\section{xtvutils.c File Reference}
\label{xtvutils_8c}\index{xtvutils.c@{xtvutils.c}}
X11 Windows image routines. 

{\tt \#include \char`\"{}xtv.h\char`\"{}}\par
\subsection*{Defines}
\begin{CompactItemize}
\item 
\#define {\bf ALL\_\-X\_\-EVENTS}\ ($\sim$0)\label{xtvutils_8c_a0}

\end{CompactItemize}
\subsection*{Functions}
\begin{CompactItemize}
\item 
int {\bf imageinit} (int $\ast$winit, int $\ast$hinit, int $\ast$nclrs, int zfinit, int yupinit, char $\ast$resourcename, char $\ast$windowname, int xoff, int yoff)
\begin{CompactList}\small\item\em Initialize the Image. \item\end{CompactList}\item 
int {\bf createwindows} (char $\ast$resourcename, char $\ast$windowname, int xoff, int yoff)
\begin{CompactList}\small\item\em Create the base window and all its little friends. \item\end{CompactList}\item 
void {\bf imageclose} ()\label{xtvutils_8c_a167}

\begin{CompactList}\small\item\em close the image window \item\end{CompactList}\item 
void {\bf imageerase} ()\label{xtvutils_8c_a168}

\begin{CompactList}\small\item\em erase the displayed image \item\end{CompactList}\item 
int {\bf imageupdate} (int x0, int y0, int nx, int ny, int map)
\begin{CompactList}\small\item\em update the image (or a subset) on the display \item\end{CompactList}\item 
int {\bf imagedisplay} (int x0, int y0, int nx, int ny, int awidth, int xoff, int yoff, float $\ast$a, int color)
\begin{CompactList}\small\item\em Display an image. \item\end{CompactList}\item 
void {\bf imagemap} (float zero, float span, int {\bf ncolors})
\begin{CompactList}\small\item\em Create the color loolup table. \item\end{CompactList}\item 
void {\bf imagepalette} (int n, short $\ast$r, short $\ast$g, short $\ast$b, int flag)
\begin{CompactList}\small\item\em Load the X Windows color palette. \item\end{CompactList}\item 
void {\bf imageinstallkey} (int key, int echo, void($\ast$function)(int, int, int, int, int))
\begin{CompactList}\small\item\em Install the image cursor action keys. \item\end{CompactList}\item 
void {\bf imageuninstallkey} (int key)\label{xtvutils_8c_a174}

\begin{CompactList}\small\item\em Uninstall an image cursor action key. \item\end{CompactList}\item 
void {\bf imagetext} (int xuser, int yuser, char $\ast$text, int $\ast$textlen)
\begin{CompactList}\small\item\em imagetext - draw text over the image \item\end{CompactList}\item 
void {\bf imagebox} (int x, int y, int nx, int ny, int color)
\begin{CompactList}\small\item\em Draw a box on the image. \item\end{CompactList}\item 
void {\bf imagecross} (int x, int y, int r)
\begin{CompactList}\small\item\em Draw a cross (+) on the image. \item\end{CompactList}\item 
void {\bf imagerelocate} (int xuser, int yuser)
\begin{CompactList}\small\item\em Move the overlay vector pointer to user X,Y coordinates. \item\end{CompactList}\item 
void {\bf imagedraw} (int xuser, int yuser, int color)
\begin{CompactList}\small\item\em Draw to user coordinates (X,Y) from the last position. \item\end{CompactList}\item 
void {\bf imagedrawflush} (int xuser, int yuser, int color)
\begin{CompactList}\small\item\em Draw and immediately refresh the display. \item\end{CompactList}\item 
void {\bf vecclear} ()
\begin{CompactList}\small\item\em Clear the overlay plane. \item\end{CompactList}\item 
void {\bf storeclear} ()
\begin{CompactList}\small\item\em clears the image storage (ring) buffer \item\end{CompactList}\item 
void {\bf imagevnull} ()\label{xtvutils_8c_a183}

\begin{CompactList}\small\item\em Clear (null) the vector and text arrays. \item\end{CompactList}\item 
void {\bf imagetreplay} ()\label{xtvutils_8c_a184}

\begin{CompactList}\small\item\em redraw text on the overlay plane \item\end{CompactList}\item 
void {\bf imagevreplay} ()\label{xtvutils_8c_a185}

\begin{CompactList}\small\item\em Redraw vectors on the overlay plane. \item\end{CompactList}\item 
void {\bf lights} (int state)
\begin{CompactList}\small\item\em Work the various status \char`\"{}lights\char`\"{} on the display. \item\end{CompactList}\item 
void {\bf imagelight} (int upper, char $\ast$mess, int color)
\begin{CompactList}\small\item\em Draw an status light on the display. \item\end{CompactList}\item 
int {\bf imageread} (int $\ast$x, int $\ast$y, char $\ast$key)
\begin{CompactList}\small\item\em Blocking read from the display of one character. \item\end{CompactList}\item 
void {\bf updateimage} (int x0, int y0, int nx, int ny, int awidth, float $\ast$a, int x1, int y1)
\begin{CompactList}\small\item\em Update the image and window dimension database. \item\end{CompactList}\item 
void {\bf updatestore} ()\label{xtvutils_8c_a190}

\begin{CompactList}\small\item\em extract image information from the ring-buffer store \item\end{CompactList}\item 
void {\bf updatepan} (int xc, int yc, int in)
\begin{CompactList}\small\item\em Update the centering of the image (pan-and-zoom). \item\end{CompactList}\item 
void {\bf updatesize} (int wid, int hgt)
\begin{CompactList}\small\item\em Update all the windows to a new overall size. \item\end{CompactList}\item 
void {\bf newsizesubwin} (int wid, int hgt)
\begin{CompactList}\small\item\em Resize subwindows. \item\end{CompactList}\item 
void {\bf resizesubwin} ()
\begin{CompactList}\small\item\em resize the subwindow \item\end{CompactList}\item 
void {\bf updatename} (char $\ast$text, int num)
\begin{CompactList}\small\item\em update the name that appears on the image display \item\end{CompactList}\item 
void {\bf tvimnum} (char $\ast$str, int n)
\begin{CompactList}\small\item\em Update the image number of the displayed image. \item\end{CompactList}\item 
void {\bf updatecoords} (int wx, int wy, int key)
\begin{CompactList}\small\item\em Update the cursor x,y coords on the status window. \item\end{CompactList}\item 
void {\bf updatebrkpt} (int x, int ibut)
\begin{CompactList}\small\item\em Display the color palette break point vector value on cursor action. \item\end{CompactList}\item 
void {\bf brkwrite} (char $\ast$s, float b, int left)
\begin{CompactList}\small\item\em Fit a floating point number into a 7-character string. \item\end{CompactList}\item 
void {\bf resetcolors} ()
\begin{CompactList}\small\item\em Reset the color palette. \item\end{CompactList}\item 
void {\bf newcolors} (int n1, int n2)
\begin{CompactList}\small\item\em Compress the color palette. \item\end{CompactList}\item 
void {\bf updatenewpal} (int x0, int ibut)
\begin{CompactList}\small\item\em update the color palette after modification \item\end{CompactList}\item 
void {\bf updatepal} (int x, int ibut)
\begin{CompactList}\small\item\em Update the palette. \item\end{CompactList}\item 
int {\bf updatezoom} (int xw, int yw)\label{xtvutils_8c_a204}

\item 
void {\bf oldzcursor} (int j)
\begin{CompactList}\small\item\em Old zoom-window cursor position (updatezoom() utility function). \item\end{CompactList}\item 
void {\bf zcursor} (int j)
\begin{CompactList}\small\item\em New zoom cursor position (updatezoom() utility function). \item\end{CompactList}\item 
void {\bf writepix} (int x, int y, int wid, int hgt)
\begin{CompactList}\small\item\em Write a portion of the image to the image window. \item\end{CompactList}\item 
int {\bf xtv\_\-refresh} (int signo)
\begin{CompactList}\small\item\em Interrupt handler for X Events. \item\end{CompactList}\item 
void {\bf keyzoompan} (int x, int y, int inout)
\begin{CompactList}\small\item\em Zoom/Pan Key Action Function. \item\end{CompactList}\item 
void {\bf keyzoomin} (int x, int y, int xuser, int yuser, int key)
\begin{CompactList}\small\item\em Zoom In Action Function. \item\end{CompactList}\item 
void {\bf keyzoomout} (int x, int y, int xuser, int yuser, int key)
\begin{CompactList}\small\item\em Zoom Out Action Function. \item\end{CompactList}\item 
void {\bf keypan} (int x, int y, int xuser, int yuser, int key)
\begin{CompactList}\small\item\em Image Pan Action Function. \item\end{CompactList}\item 
void {\bf keyrecenter} (int x, int y, int xuser, int yuser, int key)
\begin{CompactList}\small\item\em Restore Original Zoom/Center Action Function. \item\end{CompactList}\item 
void {\bf keyzoomprint} (int x, int y, int xuser, int yuser, int key)
\begin{CompactList}\small\item\em Power Zoom/Pan Action Function. \item\end{CompactList}\item 
void {\bf keyhelp} (int x, int y, int xuser, int yuser, int key)
\begin{CompactList}\small\item\em Help Action Function. \item\end{CompactList}\item 
void {\bf zfreeze} (int x, int y, int xuser, int yuser, int key)
\begin{CompactList}\small\item\em Freeze Zoom Action Function. \item\end{CompactList}\item 
void {\bf nextim} (int x, int y, int xuser, int yuser, int key)
\begin{CompactList}\small\item\em Blink to Next Image Action Function. \item\end{CompactList}\item 
void {\bf lastim} (int x, int y, int xuser, int yuser, int key)
\begin{CompactList}\small\item\em Blink to Previous Image Action Function. \item\end{CompactList}\item 
void {\bf zpeak} (int x, int y, int xuser, int yuser, int key)
\begin{CompactList}\small\item\em Find Peak Pixel Nearest the Cursor Action Function. \item\end{CompactList}\item 
void {\bf zhairs} (int x, int y, int xuser, int yuser, int key)
\begin{CompactList}\small\item\em Toggle Full-Screen Cross Hairs On/Off Action Function. \item\end{CompactList}\item 
void {\bf vnohair} ()\label{xtvutils_8c_a221}

\begin{CompactList}\small\item\em ... \item\end{CompactList}\item 
int {\bf black\-Pixel} (Display $\ast${\bf dpy}, int {\bf screen})
\begin{CompactList}\small\item\em Background pixmap (black). \item\end{CompactList}\item 
int {\bf white\-Pixel} (Display $\ast${\bf dpy}, int {\bf screen})
\begin{CompactList}\small\item\em Border pixmap (white). \item\end{CompactList}\item 
void {\bf mapimage} (int x0, int y0, int nx, int ny, int datw, float $\ast${\bf data}, int imx0, int imy0, int imw, int yinv, unsigned long $\ast${\bf image}, float $\ast${\bf breakpts}, unsigned long nbreak, unsigned long $\ast${\bf pixels}, int bits\_\-per\_\-pixel, unsigned long cmask)
\begin{CompactList}\small\item\em Map floating image data into integer display data. \item\end{CompactList}\item 
void {\bf replicate} (int x0, int y0, int w0, int h0, char $\ast$from, int f, int x, int y, int w, int h, int wto, char $\ast$to, int fill, int bits\_\-per\_\-pixel)
\begin{CompactList}\small\item\em replicate \item\end{CompactList}\item 
void {\bf samplicate} (int x0, int y0, int w0, int h0, char $\ast$from, int f, int x, int y, int w, int h, int wto, char $\ast$to, int fill, int bits\_\-per\_\-pixel)
\begin{CompactList}\small\item\em samplicate an array (sample instead of replicate) \item\end{CompactList}\item 
void {\bf duplicate} (int x0, int y0, int w0, int h0, char $\ast$from, int x, int y, int w, int h, int wto, char $\ast$to, int fill, int bits\_\-per\_\-pixel)
\begin{CompactList}\small\item\em duplicate pixel replication mapping without expansion/compression \item\end{CompactList}\item 
void {\bf zeroborder} (int x0, int y0, int w0, int h0, int w, int h, int wto, char $\ast$to, int nbytes)
\begin{CompactList}\small\item\em fill border with zeros \item\end{CompactList}\item 
void {\bf locate} (float xx[$\,$], unsigned long n, float x, unsigned long $\ast$j)\label{xtvutils_8c_a229}

\begin{CompactList}\small\item\em locate index of x in array xx xx data vector n length of data vector xx x value in xx to find j index of x in xx \item\end{CompactList}\item 
void {\bf hunt} (float xx[$\,$], unsigned long n, float x, unsigned long $\ast$jlo)
\begin{CompactList}\small\item\em hunt for nearest left index of x in array xx \item\end{CompactList}\item 
int {\bf xy2index} (int xuser, int yuser)
\begin{CompactList}\small\item\em Convert (x,y) user coordinates into data vector index. \item\end{CompactList}\end{CompactItemize}
\subsection*{Variables}
\begin{CompactItemize}
\item 
int {\bf usehairs} = 0\label{xtvutils_8c_a3}

\item 
int {\bf privcmap} = 0\label{xtvutils_8c_a4}

\item 
Display $\ast$ {\bf dpy}\label{xtvutils_8c_a5}

\begin{CompactList}\small\item\em Display pointer. \item\end{CompactList}\item 
Window {\bf wbase}\label{xtvutils_8c_a6}

\begin{CompactList}\small\item\em Base window. \item\end{CompactList}\item 
Window {\bf wimage}\label{xtvutils_8c_a7}

\begin{CompactList}\small\item\em Image window. \item\end{CompactList}\item 
Window {\bf wpal}\label{xtvutils_8c_a8}

\begin{CompactList}\small\item\em color palette windowd \item\end{CompactList}\item 
Window {\bf wxyz}\label{xtvutils_8c_a9}

\begin{CompactList}\small\item\em xyz data display windows \item\end{CompactList}\item 
Window {\bf wzoom}\label{xtvutils_8c_a10}

\begin{CompactList}\small\item\em Zoom window. \item\end{CompactList}\item 
Window {\bf wlgt1}\label{xtvutils_8c_a11}

\begin{CompactList}\small\item\em Status Light 1 window. \item\end{CompactList}\item 
Window {\bf wlgt2}\label{xtvutils_8c_a12}

\begin{CompactList}\small\item\em Status Light 2 window. \item\end{CompactList}\item 
Window {\bf wlgt3}\label{xtvutils_8c_a13}

\begin{CompactList}\small\item\em Status Light 3 window. \item\end{CompactList}\item 
Window {\bf wlgt4}\label{xtvutils_8c_a14}

\begin{CompactList}\small\item\em Status Light 4 window. \item\end{CompactList}\item 
int {\bf screen}\label{xtvutils_8c_a15}

\begin{CompactList}\small\item\em Screen number of display. \item\end{CompactList}\item 
XImage $\ast$ {\bf dataimage}\label{xtvutils_8c_a16}

\begin{CompactList}\small\item\em data image \item\end{CompactList}\item 
XImage $\ast$ {\bf palimage}\label{xtvutils_8c_a17}

\begin{CompactList}\small\item\em palette strip image structure \item\end{CompactList}\item 
XImage $\ast$ {\bf zoomimage}\label{xtvutils_8c_a18}

\begin{CompactList}\small\item\em zoom image \item\end{CompactList}\item 
GC {\bf vectorgc}\label{xtvutils_8c_a19}

\item 
GC {\bf textgc}\label{xtvutils_8c_a20}

\begin{CompactList}\small\item\em Graphics context of the vector and text planes. \item\end{CompactList}\item 
GC {\bf imagegc}\label{xtvutils_8c_a21}

\begin{CompactList}\small\item\em Image graphics context of the image. \item\end{CompactList}\item 
GC {\bf xorgc}\label{xtvutils_8c_a22}

\begin{CompactList}\small\item\em Image graphics context of the cross-hairs. \item\end{CompactList}\item 
Font {\bf font}\label{xtvutils_8c_a23}

\item 
char $\ast$ {\bf fontname1}\label{xtvutils_8c_a24}

\item 
char $\ast$ {\bf fontname2}\label{xtvutils_8c_a25}

\item 
char $\ast$$\ast$ {\bf fontpath}\label{xtvutils_8c_a26}

\item 
int $\ast$ {\bf nfont}\label{xtvutils_8c_a27}

\item 
XFont\-Struct $\ast$ {\bf fontinfo}\label{xtvutils_8c_a28}

\begin{CompactList}\small\item\em font info struct \item\end{CompactList}\item 
int {\bf fontheight}\label{xtvutils_8c_a29}

\begin{CompactList}\small\item\em font height \item\end{CompactList}\item 
int {\bf fontwidth}\label{xtvutils_8c_a30}

\begin{CompactList}\small\item\em font width \item\end{CompactList}\item 
int {\bf ncolors}\label{xtvutils_8c_a31}

\begin{CompactList}\small\item\em number of image colors used \item\end{CompactList}\item 
int {\bf nvcolors}\label{xtvutils_8c_a32}

\begin{CompactList}\small\item\em number of vector colors used \item\end{CompactList}\item 
unsigned long {\bf pixels} [256]\label{xtvutils_8c_a33}

\begin{CompactList}\small\item\em color cell addresses allocated \item\end{CompactList}\item 
unsigned long {\bf planes} [1]\label{xtvutils_8c_a34}

\begin{CompactList}\small\item\em plane mask \item\end{CompactList}\item 
Visual $\ast$ {\bf visual}\label{xtvutils_8c_a35}

\begin{CompactList}\small\item\em window visual (e.g., Pseudo\-Color, True\-Color, etc) \item\end{CompactList}\item 
int {\bf depth}\label{xtvutils_8c_a36}

\begin{CompactList}\small\item\em window bit depth (e.g., 16- or 24-bit) \item\end{CompactList}\item 
Colormap {\bf defcmap}\label{xtvutils_8c_a37}

\begin{CompactList}\small\item\em Display color map. \item\end{CompactList}\item 
XColor {\bf cmap} [256]\label{xtvutils_8c_a38}

\begin{CompactList}\small\item\em color cells used \item\end{CompactList}\item 
XColor {\bf stcolor} [5]\label{xtvutils_8c_a39}

\begin{CompactList}\small\item\em standard colors status lights, overlay \item\end{CompactList}\item 
XColor {\bf palcolors} [MAXCOLORS]\label{xtvutils_8c_a40}

\begin{CompactList}\small\item\em complete palette \item\end{CompactList}\item 
XColor {\bf vcolor} [256]\label{xtvutils_8c_a41}

\begin{CompactList}\small\item\em vector color palette \item\end{CompactList}\item 
int {\bf npalette}\label{xtvutils_8c_a42}

\begin{CompactList}\small\item\em size of complete palette \item\end{CompactList}\item 
int {\bf pal0}\label{xtvutils_8c_a43}

\item 
int {\bf pal1}\label{xtvutils_8c_a44}

\begin{CompactList}\small\item\em palette breakpoints \item\end{CompactList}\item 
float {\bf breakpts} [256]\label{xtvutils_8c_a45}

\begin{CompactList}\small\item\em break points for the color palette lookup table \item\end{CompactList}\item 
int {\bf truecolor}\label{xtvutils_8c_a46}

\begin{CompactList}\small\item\em Flag: =1 if a True\-Color visual. \item\end{CompactList}\item 
int {\bf directcolor}\label{xtvutils_8c_a47}

\begin{CompactList}\small\item\em Flag: =1 if a Direct\-Color visual. \item\end{CompactList}\item 
unsigned long {\bf rmask}\label{xtvutils_8c_a48}

\begin{CompactList}\small\item\em Red color plane mask (from XVisual\-Info struct). \item\end{CompactList}\item 
unsigned long {\bf gmask}\label{xtvutils_8c_a49}

\begin{CompactList}\small\item\em Green color plane mask (from XVisual\-Info struct). \item\end{CompactList}\item 
unsigned long {\bf bmask}\label{xtvutils_8c_a50}

\begin{CompactList}\small\item\em Blue color plane mask (from XVisual\-Info struct). \item\end{CompactList}\item 
XError\-Event {\bf \_\-XError\-Event}\label{xtvutils_8c_a51}

\begin{CompactList}\small\item\em X Error Event struct. \item\end{CompactList}\item 
{\bf imagevector\_\-t} {\bf veclist} [MAXVEC]\label{xtvutils_8c_a52}

\begin{CompactList}\small\item\em Save array for vectors. \item\end{CompactList}\item 
int {\bf xvlast}\label{xtvutils_8c_a53}

\item 
int {\bf yvlast}\label{xtvutils_8c_a54}

\item 
int {\bf vcount}\label{xtvutils_8c_a55}

\begin{CompactList}\small\item\em x,y of last vector; number drawn \item\end{CompactList}\item 
int {\bf vlastcolor}\label{xtvutils_8c_a56}

\begin{CompactList}\small\item\em color of last vector \item\end{CompactList}\item 
{\bf imagetext\_\-t} {\bf textlist} [MAXTEXT]\label{xtvutils_8c_a57}

\begin{CompactList}\small\item\em Save array for text. \item\end{CompactList}\item 
int {\bf xtlast}\label{xtvutils_8c_a58}

\item 
int {\bf ytlast}\label{xtvutils_8c_a59}

\item 
int {\bf tcount}\label{xtvutils_8c_a60}

\begin{CompactList}\small\item\em x,y of last text; number drawn \item\end{CompactList}\item 
int {\bf tlastcolor}\label{xtvutils_8c_a61}

\begin{CompactList}\small\item\em color of last text \item\end{CompactList}\item 
int {\bf maxwidth}\label{xtvutils_8c_a62}

\item 
int {\bf maxheight}\label{xtvutils_8c_a63}

\begin{CompactList}\small\item\em maximum image size \item\end{CompactList}\item 
int {\bf width}\label{xtvutils_8c_a64}

\item 
int {\bf height}\label{xtvutils_8c_a65}

\begin{CompactList}\small\item\em last size of image part of window \item\end{CompactList}\item 
int {\bf resize}\label{xtvutils_8c_a66}

\begin{CompactList}\small\item\em flag for auto resizing of window \item\end{CompactList}\item 
int {\bf autozoomout}\label{xtvutils_8c_a67}

\begin{CompactList}\small\item\em flag for autozoomout of large images \item\end{CompactList}\item 
int {\bf zoomsample}\label{xtvutils_8c_a68}

\begin{CompactList}\small\item\em flag for sample/bin of large images \item\end{CompactList}\item 
int {\bf to\_\-program}\label{xtvutils_8c_a69}

\item 
int {\bf from\_\-display}\label{xtvutils_8c_a70}

\begin{CompactList}\small\item\em display comm pipe descriptors \item\end{CompactList}\item 
int {\bf waiting\_\-for\_\-key} = 0\label{xtvutils_8c_a71}

\begin{CompactList}\small\item\em flag for blocking read \item\end{CompactList}\item 
int {\bf store}\label{xtvutils_8c_a72}

\begin{CompactList}\small\item\em storage counter, runs 0..{\bf MAXSTORE}{\rm (p.\,\pageref{xtv_8h_a28})} \item\end{CompactList}\item 
float $\ast$ {\bf data}\label{xtvutils_8c_a73}

\begin{CompactList}\small\item\em address of data (image) array \item\end{CompactList}\item 
char $\ast$ {\bf image} [MAXSTORE]\label{xtvutils_8c_a74}

\begin{CompactList}\small\item\em data converted to display bytes \item\end{CompactList}\item 
int {\bf nimage} [MAXSTORE]\label{xtvutils_8c_a75}

\begin{CompactList}\small\item\em number of bytes in each image \item\end{CompactList}\item 
char $\ast$ {\bf imbuf}\label{xtvutils_8c_a76}

\item 
char $\ast$ {\bf outbuf}\label{xtvutils_8c_a77}

\begin{CompactList}\small\item\em 128k buffer for display \item\end{CompactList}\item 
int {\bf yup}\label{xtvutils_8c_a78}

\begin{CompactList}\small\item\em flag = 0/1 for y increasing down/up \item\end{CompactList}\item 
int {\bf datw}\label{xtvutils_8c_a79}

\item 
int {\bf dath}\label{xtvutils_8c_a80}

\begin{CompactList}\small\item\em dimensions of data array \item\end{CompactList}\item 
int {\bf imwidth}\label{xtvutils_8c_a81}

\begin{CompactList}\small\item\em bytes per row of image array=4$\ast$((imw+3)/4) \item\end{CompactList}\item 
int {\bf imw}\label{xtvutils_8c_a82}

\item 
int {\bf imh}\label{xtvutils_8c_a83}

\begin{CompactList}\small\item\em dimensions of image array \item\end{CompactList}\item 
int {\bf winw}\label{xtvutils_8c_a84}

\item 
int {\bf winh}\label{xtvutils_8c_a85}

\begin{CompactList}\small\item\em dimensions of image in window \item\end{CompactList}\item 
int {\bf zim}\label{xtvutils_8c_a86}

\begin{CompactList}\small\item\em expansion factor of image \item\end{CompactList}\item 
int {\bf offx}\label{xtvutils_8c_a87}

\item 
int {\bf offy}\label{xtvutils_8c_a88}

\begin{CompactList}\small\item\em original coords of UL of data array \item\end{CompactList}\item 
int {\bf datx}\label{xtvutils_8c_a89}

\item 
int {\bf daty}\label{xtvutils_8c_a90}

\begin{CompactList}\small\item\em data coords of UL of image portion \item\end{CompactList}\item 
int {\bf imx}\label{xtvutils_8c_a91}

\item 
int {\bf imy}\label{xtvutils_8c_a92}

\begin{CompactList}\small\item\em image coords of UL of displayed portion \item\end{CompactList}\item 
int {\bf winx}\label{xtvutils_8c_a93}

\item 
int {\bf winy}\label{xtvutils_8c_a94}

\begin{CompactList}\small\item\em window coords of UL of displayed portion \item\end{CompactList}\item 
int {\bf imtv}\label{xtvutils_8c_a95}

\begin{CompactList}\small\item\em VISTA buffer number. \item\end{CompactList}\item 
float $\ast$ {\bf sdata} [MAXSTORE]\label{xtvutils_8c_a96}

\item 
int {\bf sdatw} [MAXSTORE]\label{xtvutils_8c_a97}

\item 
int {\bf sdath} [MAXSTORE]\label{xtvutils_8c_a98}

\item 
int {\bf simwidth} [MAXSTORE]\label{xtvutils_8c_a99}

\item 
int {\bf simw} [MAXSTORE]\label{xtvutils_8c_a100}

\item 
int {\bf simh} [MAXSTORE]\label{xtvutils_8c_a101}

\item 
int {\bf swinw} [MAXSTORE]\label{xtvutils_8c_a102}

\item 
int {\bf swinh} [MAXSTORE]\label{xtvutils_8c_a103}

\item 
int {\bf szim} [MAXSTORE]\label{xtvutils_8c_a104}

\item 
int {\bf soffx} [MAXSTORE]\label{xtvutils_8c_a105}

\item 
int {\bf soffy} [MAXSTORE]\label{xtvutils_8c_a106}

\item 
int {\bf sdatx} [MAXSTORE]\label{xtvutils_8c_a107}

\item 
int {\bf sdaty} [MAXSTORE]\label{xtvutils_8c_a108}

\item 
int {\bf simx} [MAXSTORE]\label{xtvutils_8c_a109}

\item 
int {\bf simy} [MAXSTORE]\label{xtvutils_8c_a110}

\item 
int {\bf swinx} [MAXSTORE]\label{xtvutils_8c_a111}

\item 
int {\bf swiny} [MAXSTORE]\label{xtvutils_8c_a112}

\item 
int {\bf simtv} [MAXSTORE]\label{xtvutils_8c_a113}

\item 
char {\bf imname1} [MAXSTORE][MAXSTRINGWID]\label{xtvutils_8c_a114}

\item 
char {\bf imname2} [MAXSTORE][MAXSTRINGWID]\label{xtvutils_8c_a115}

\item 
int {\bf palheight}\label{xtvutils_8c_a116}

\item 
int {\bf palwidth}\label{xtvutils_8c_a117}

\item 
int {\bf palx}\label{xtvutils_8c_a118}

\item 
int {\bf paly}\label{xtvutils_8c_a119}

\begin{CompactList}\small\item\em size and loc of palette subwindow \item\end{CompactList}\item 
float $\ast$ {\bf palbrkpt}\label{xtvutils_8c_a120}

\begin{CompactList}\small\item\em pointer to breakpoint array \item\end{CompactList}\item 
char $\ast$ {\bf palette}\label{xtvutils_8c_a121}

\begin{CompactList}\small\item\em array with palette values \item\end{CompactList}\item 
short $\ast$ {\bf lookup}\label{xtvutils_8c_a122}

\begin{CompactList}\small\item\em lookup table: data -$>$ color addr \item\end{CompactList}\item 
int {\bf lgtheight}\label{xtvutils_8c_a123}

\item 
int {\bf lgtwidth}\label{xtvutils_8c_a124}

\item 
int {\bf lgtx}\label{xtvutils_8c_a125}

\item 
int {\bf lgty} [4]\label{xtvutils_8c_a126}

\begin{CompactList}\small\item\em size and loc of light subwindows \item\end{CompactList}\item 
char $\ast$ {\bf uppermess1} = NULL\label{xtvutils_8c_a128}

\item 
char $\ast$ {\bf lowermess1} = NULL\label{xtvutils_8c_a129}

\item 
char $\ast$ {\bf uppermess2} = NULL\label{xtvutils_8c_a130}

\item 
char $\ast$ {\bf lowermess2} = NULL\label{xtvutils_8c_a131}

\item 
char {\bf xyzstring} [NXYZ]\label{xtvutils_8c_a132}

\item 
char {\bf numstr} [20]\label{xtvutils_8c_a133}

\begin{CompactList}\small\item\em strings used for the cursor data display \item\end{CompactList}\item 
int {\bf xyzheight}\label{xtvutils_8c_a134}

\item 
int {\bf xyzwidth}\label{xtvutils_8c_a135}

\item 
int {\bf xyzx}\label{xtvutils_8c_a136}

\item 
int {\bf xyzy}\label{xtvutils_8c_a137}

\begin{CompactList}\small\item\em size and loc of xyz subwindow \item\end{CompactList}\item 
int {\bf mousex}\label{xtvutils_8c_a138}

\item 
int {\bf mousey}\label{xtvutils_8c_a139}

\item 
int {\bf mousez}\label{xtvutils_8c_a140}

\begin{CompactList}\small\item\em loc and value of mouse \item\end{CompactList}\item 
int {\bf textx}\label{xtvutils_8c_a141}

\item 
int {\bf texty} [3]\label{xtvutils_8c_a142}

\begin{CompactList}\small\item\em loc of xyz text \item\end{CompactList}\item 
char $\ast$ {\bf zimage}\label{xtvutils_8c_a143}

\begin{CompactList}\small\item\em zoom buffer \item\end{CompactList}\item 
int {\bf zoomf}\label{xtvutils_8c_a144}

\begin{CompactList}\small\item\em zoom factor \item\end{CompactList}\item 
int {\bf zwidth}\label{xtvutils_8c_a145}

\item 
int {\bf zheight}\label{xtvutils_8c_a146}

\begin{CompactList}\small\item\em zoom window size \item\end{CompactList}\item 
int {\bf lastzoomx}\label{xtvutils_8c_a147}

\item 
int {\bf lastzoomy}\label{xtvutils_8c_a148}

\begin{CompactList}\small\item\em location of last update of zoom \item\end{CompactList}\item 
{\bf imagevector\_\-t} {\bf zcursvec} [MAXZCURS]\label{xtvutils_8c_a149}

\begin{CompactList}\small\item\em array of zoom cursor vectors \item\end{CompactList}\item 
int {\bf whichkey}\label{xtvutils_8c_a150}

\begin{CompactList}\small\item\em which button/key was pressed \item\end{CompactList}\item 
int {\bf buttondown}\label{xtvutils_8c_a151}

\begin{CompactList}\small\item\em true if a button is pressed \item\end{CompactList}\item 
int {\bf lastx}\label{xtvutils_8c_a152}

\item 
int {\bf lasty}\label{xtvutils_8c_a153}

\begin{CompactList}\small\item\em loc where button/key went down \item\end{CompactList}\item 
{\bf keyaction\_\-t} {\bf keyaction} [128]\label{xtvutils_8c_a154}

\begin{CompactList}\small\item\em Action keys. \item\end{CompactList}\item 
int {\bf TVis\-Open} = 0\label{xtvutils_8c_a155}

\item 
int {\bf fd\_\-for\_\-X} = -1\label{xtvutils_8c_a156}

\item 
int {\bf imagevalid}\label{xtvutils_8c_a157}

\item 
XColor {\bf testcolor} [200]\label{xtvutils_8c_a158}

\item 
int {\bf tvinit} = 0\label{xtvutils_8c_a159}

\item 
int {\bf zfreezeon} = 0\label{xtvutils_8c_a160}

\item 
int {\bf update} = 0\label{xtvutils_8c_a161}

\item 
int {\bf id} = 0\label{xtvutils_8c_a163}

\end{CompactItemize}


\subsection{Detailed Description}
X11 Windows image routines. 

\begin{Desc}
\item[Author:]John Tonry 5/12/88\end{Desc}
Implemented in xvista with various mods - J. Holtzman 5/90-5/91, 96

Extracted into standalone version and added doxygen markup\par
 R. Pogge, OSU Astronomy Dept. ({\tt pogge@astronomy.ohio.state.edu})\par
 2005 May 27

\subsection{Function Documentation}
\index{xtvutils.c@{xtvutils.c}!blackPixel@{blackPixel}}
\index{blackPixel@{blackPixel}!xtvutils.c@{xtvutils.c}}
\subsubsection{\setlength{\rightskip}{0pt plus 5cm}int black\-Pixel (Display $\ast$ {\em dpy}, int {\em screen})}\label{xtvutils_8c_a222}


Background pixmap (black). 

\begin{Desc}
\item[Parameters:]
\begin{description}
\item[{\em dpy}]Pointer to the display \item[{\em screen}]screen number \end{description}
\end{Desc}
\index{xtvutils.c@{xtvutils.c}!brkwrite@{brkwrite}}
\index{brkwrite@{brkwrite}!xtvutils.c@{xtvutils.c}}
\subsubsection{\setlength{\rightskip}{0pt plus 5cm}void brkwrite (char $\ast$ {\em s}, float {\em b}, int {\em left})}\label{xtvutils_8c_a199}


Fit a floating point number into a 7-character string. 

\begin{Desc}
\item[Parameters:]
\begin{description}
\item[{\em s}]character string to create \item[{\em b}]floating number to pack into 7 digits \item[{\em left}]1=left pad with blanks\end{description}
\end{Desc}
Converts a floating number (b) into a 7-character string with the appropriate format to not overrun the string size. The optional left flag is set if the string is to be padded on the left with spaces. \index{xtvutils.c@{xtvutils.c}!createwindows@{createwindows}}
\index{createwindows@{createwindows}!xtvutils.c@{xtvutils.c}}
\subsubsection{\setlength{\rightskip}{0pt plus 5cm}int createwindows (char $\ast$ {\em resourcename}, char $\ast$ {\em windowname}, int {\em xoff}, int {\em yoff})}\label{xtvutils_8c_a166}


Create the base window and all its little friends. 

\begin{Desc}
\item[Parameters:]
\begin{description}
\item[{\em resourcename}]Name of the X-resource file to use \item[{\em windowname}]Name of the window (appears on the top bar) \item[{\em xoff}]X-axis offset on the desktop \item[{\em yoff}]Y-axis offset on the desktop\end{description}
\end{Desc}
\begin{Desc}
\item[Returns:]0 on success, error codes on failure.\end{Desc}
Creates the base window and all the little subwindows that are needed by the image display.

$<$ set windows attributes \index{xtvutils.c@{xtvutils.c}!duplicate@{duplicate}}
\index{duplicate@{duplicate}!xtvutils.c@{xtvutils.c}}
\subsubsection{\setlength{\rightskip}{0pt plus 5cm}void duplicate (int {\em x0}, int {\em y0}, int {\em w0}, int {\em h0}, char $\ast$ {\em from}, int {\em x}, int {\em y}, int {\em w}, int {\em h}, int {\em wto}, char $\ast$ {\em to}, int {\em fill}, int {\em bits\_\-per\_\-pixel})}\label{xtvutils_8c_a227}


duplicate pixel replication mapping without expansion/compression 

\begin{Desc}
\item[Parameters:]
\begin{description}
\item[{\em x0}]X-axis origin of the source image \item[{\em y0}]Y-axis origin of the source image \item[{\em w0}]width (X-axis size) of source image \item[{\em h0}]height (Y-axis size) of source image \item[{\em from}]Source image array \item[{\em x}]X-axis origin of the destination image \item[{\em y}]Y-axis origin of the destination image \item[{\em w}]width (X-axis size) of the destination image \item[{\em h}]height (Y-axis size) of the destination image \item[{\em wto}]dimensions of the destination image \item[{\em to}]Destination image array \item[{\em fill}]Flag: 1=zero rest of the to array \item[{\em bits\_\-per\_\-pixel}]data precision of the arrays\end{description}
\end{Desc}
The same as replicate with no expansion or compression \index{xtvutils.c@{xtvutils.c}!hunt@{hunt}}
\index{hunt@{hunt}!xtvutils.c@{xtvutils.c}}
\subsubsection{\setlength{\rightskip}{0pt plus 5cm}void hunt (float {\em xx}[$\,$], unsigned long {\em n}, float {\em x}, unsigned long $\ast$ {\em jlo})}\label{xtvutils_8c_a230}


hunt for nearest left index of x in array xx 

\begin{Desc}
\item[Parameters:]
\begin{description}
\item[{\em xx}]data vector \item[{\em n}]length of the data vector xx \item[{\em x}]data value to find in the vector \item[{\em jlo}](returned) index of x \end{description}
\end{Desc}
\index{xtvutils.c@{xtvutils.c}!imagebox@{imagebox}}
\index{imagebox@{imagebox}!xtvutils.c@{xtvutils.c}}
\subsubsection{\setlength{\rightskip}{0pt plus 5cm}void imagebox (int {\em x}, int {\em y}, int {\em nx}, int {\em ny}, int {\em color})}\label{xtvutils_8c_a176}


Draw a box on the image. 

\begin{Desc}
\item[Parameters:]
\begin{description}
\item[{\em x}]X-axis coordinates of the lower corner of the box \item[{\em y}]Y-axis coordinates of the lower corner of the box \item[{\em nx}]width of the box in X \item[{\em ny}]width of the box in Y \item[{\em color}]color (palette index) to draw\end{description}
\end{Desc}
Draw a box on the image (vector plane) with the indicated location, size and color. \index{xtvutils.c@{xtvutils.c}!imagecross@{imagecross}}
\index{imagecross@{imagecross}!xtvutils.c@{xtvutils.c}}
\subsubsection{\setlength{\rightskip}{0pt plus 5cm}void imagecross (int {\em x}, int {\em y}, int {\em r})}\label{xtvutils_8c_a177}


Draw a cross (+) on the image. 

\begin{Desc}
\item[Parameters:]
\begin{description}
\item[{\em x}]X coordinates of the cross center \item[{\em y}]Y coordinates of the cross center \item[{\em r}]Radius of the cross (pixels)\end{description}
\end{Desc}
Draws a cross (+) on the image (vector plane) centered at the given (x,y) coordinates and the given radius. \index{xtvutils.c@{xtvutils.c}!imagedisplay@{imagedisplay}}
\index{imagedisplay@{imagedisplay}!xtvutils.c@{xtvutils.c}}
\subsubsection{\setlength{\rightskip}{0pt plus 5cm}int imagedisplay (int {\em x0}, int {\em y0}, int {\em nx}, int {\em ny}, int {\em awidth}, int {\em xoff}, int {\em yoff}, float $\ast$ {\em a}, int {\em color})}\label{xtvutils_8c_a170}


Display an image. 

\begin{Desc}
\item[Parameters:]
\begin{description}
\item[{\em a}]Input floating image array \item[{\em x0}]X-axis origin of the area to be displayed \item[{\em y0}]Y-axis origin of the area to be displayed \item[{\em nx}]X-axis size of the area to be displayed \item[{\em ny}]Y-axis size of area to be displayed \item[{\em awidth}]Number of pixels per row in a \item[{\em xoff}]X-axis coordinates to be used for pixel (0,0) \item[{\em yoff}]Y-axis coordinates to be used for pixel (0,0) \item[{\em color}]Color plane to load (1=R,2=G,3=B,0=all?)\end{description}
\end{Desc}
Does the dirty work of image display \index{xtvutils.c@{xtvutils.c}!imagedraw@{imagedraw}}
\index{imagedraw@{imagedraw}!xtvutils.c@{xtvutils.c}}
\subsubsection{\setlength{\rightskip}{0pt plus 5cm}void imagedraw (int {\em xuser}, int {\em yuser}, int {\em color})}\label{xtvutils_8c_a179}


Draw to user coordinates (X,Y) from the last position. 

\begin{Desc}
\item[Parameters:]
\begin{description}
\item[{\em xuser}]X-axis coordinate in user space \item[{\em yuser}]Y-axis coordinate in user space \item[{\em color}]color to use (from the palette)\end{description}
\end{Desc}
Draws a line from the last vector position (previous {\bf imagerelocate()}{\rm (p.\,\pageref{xtv_8h_a72})} or {\bf imagedraw()}{\rm (p.\,\pageref{xtv_8h_a73})} call) to the given (X,Y) location with a particular palette color.

The line will not appear immediately on the overlay plane unless the display is refreshed.

\begin{Desc}
\item[See also:]{\bf imagerelocate()}{\rm (p.\,\pageref{xtv_8h_a72})}, {\bf imagedrawflush()}{\rm (p.\,\pageref{xtv_8h_a74})} \end{Desc}
\index{xtvutils.c@{xtvutils.c}!imagedrawflush@{imagedrawflush}}
\index{imagedrawflush@{imagedrawflush}!xtvutils.c@{xtvutils.c}}
\subsubsection{\setlength{\rightskip}{0pt plus 5cm}void imagedrawflush (int {\em xuser}, int {\em yuser}, int {\em color})}\label{xtvutils_8c_a180}


Draw and immediately refresh the display. 

\begin{Desc}
\item[Parameters:]
\begin{description}
\item[{\em xuser}]X-axis coordinate in user space \item[{\em yuser}]Y-axis coordinate in user space \item[{\em color}]color to use (from the palette)\end{description}
\end{Desc}
Draws a line from the last vector position (previous {\bf imagerelocate()}{\rm (p.\,\pageref{xtv_8h_a72})} or {\bf imagedraw()}{\rm (p.\,\pageref{xtv_8h_a73})} call) to the given (X,Y) location with a particular palette color, and immediately refresh the display.

\begin{Desc}
\item[See also:]{\bf imagerelocate()}{\rm (p.\,\pageref{xtv_8h_a72})}, {\bf imagedraw()}{\rm (p.\,\pageref{xtv_8h_a73})} \end{Desc}
\index{xtvutils.c@{xtvutils.c}!imageinit@{imageinit}}
\index{imageinit@{imageinit}!xtvutils.c@{xtvutils.c}}
\subsubsection{\setlength{\rightskip}{0pt plus 5cm}int imageinit (int $\ast$ {\em winit}, int $\ast$ {\em hinit}, int $\ast$ {\em nclrs}, int {\em zfinit}, int {\em yupinit}, char $\ast$ {\em resourcename}, char $\ast$ {\em windowname}, int {\em xoff}, int {\em yoff})}\label{xtvutils_8c_a165}


Initialize the Image. 

\begin{Desc}
\item[Parameters:]
\begin{description}
\item[{\em winit}]width of the window (request on input, actual returned) \item[{\em hinit}]height of the window (request on input, actual returned) \item[{\em nclrs}]number of image color cells (requested on input, actual returned) \item[{\em zfinit}]initial zoom factor (0 for no zoom window) \item[{\em yupinit}]flag = 0/1 for y increasing down/up \item[{\em resourcename}]name of the X resource for this window \item[{\em windowname}]name of the window (to appear on the top bar) \item[{\em xoff}]X-axis offset on the desktop on startup \item[{\em yoff}]Y-axis offset on the desktop on startup\end{description}
\end{Desc}
\begin{Desc}
\item[Returns:]file descriptor for X events on success, negative error codes on failure.\end{Desc}
Initializes the window and its resources for the image display and sets up the event handler. \index{xtvutils.c@{xtvutils.c}!imageinstallkey@{imageinstallkey}}
\index{imageinstallkey@{imageinstallkey}!xtvutils.c@{xtvutils.c}}
\subsubsection{\setlength{\rightskip}{0pt plus 5cm}void imageinstallkey (int {\em key}, int {\em echo}, void($\ast$)(int, int, int, int, int) {\em function})}\label{xtvutils_8c_a173}


Install the image cursor action keys. 

\begin{Desc}
\item[Parameters:]
\begin{description}
\item[{\em key}]ASCII code of the key to bind to the action function \item[{\em echo}]flag: 0/1 to echo key in x,y,z display \item[{\em function()}]action function to be called when key is pressed.\end{description}
\end{Desc}
The function prototype is \small\begin{alltt}
  void action(int x, int y, int xuser, int yuser, int key) 
  \end{alltt}\normalsize 
 \index{xtvutils.c@{xtvutils.c}!imagelight@{imagelight}}
\index{imagelight@{imagelight}!xtvutils.c@{xtvutils.c}}
\subsubsection{\setlength{\rightskip}{0pt plus 5cm}void imagelight (int {\em upper}, char $\ast$ {\em mess}, int {\em color})}\label{xtvutils_8c_a187}


Draw an status light on the display. 

\begin{Desc}
\item[Parameters:]
\begin{description}
\item[{\em upper}]tells which light to draw in (see below) \item[{\em mess}]String to be written (null terminated) \item[{\em color}]0/1 for red/green\end{description}
\end{Desc}
Draws a status light. The values of upper give which light is to be operated: \small\begin{alltt}
  upper = -1 -$>$ redraw all
  upper = 1  -$>$ top light
  upper = 2  -$>$ next from top
  upper = 3  -$>$ next from bottom
  upper = 4  -$>$ bottom light
  \end{alltt}\normalsize 
 Or, graphically \small\begin{alltt}
  +-------+
  |   1   |
  +-------+
  |   2   |
  +-------+
  |   3   |
  +-------+
  |   4   |
  +-------+
  \end{alltt}\normalsize 
 \index{xtvutils.c@{xtvutils.c}!imagemap@{imagemap}}
\index{imagemap@{imagemap}!xtvutils.c@{xtvutils.c}}
\subsubsection{\setlength{\rightskip}{0pt plus 5cm}void imagemap (float {\em zero}, float {\em span}, int {\em ncolors})}\label{xtvutils_8c_a171}


Create the color loolup table. 

\begin{Desc}
\item[Parameters:]
\begin{description}
\item[{\em zero}]zero point (min) of the intensity mapping \item[{\em span}]span of the intensity mapping (max-min) \item[{\em ncolors}]Number of colors to map\end{description}
\end{Desc}
Creates the global breakpts array that is used to quickly map image data values into display colors. The breakpts array contains those data values that map into each \char`\"{}color\char`\"{} index in the palette vector.

The mapping is: \small\begin{alltt}
  breakpt[i] = zero + (span/(ncolors-2)) * i
  \end{alltt}\normalsize 
 Note that this algorithm reserves the top 2 colors in the palette for non-data uses. \index{xtvutils.c@{xtvutils.c}!imagepalette@{imagepalette}}
\index{imagepalette@{imagepalette}!xtvutils.c@{xtvutils.c}}
\subsubsection{\setlength{\rightskip}{0pt plus 5cm}void imagepalette (int {\em n}, short $\ast$ {\em r}, short $\ast$ {\em g}, short $\ast$ {\em b}, int {\em flag})}\label{xtvutils_8c_a172}


Load the X Windows color palette. 

\begin{Desc}
\item[Parameters:]
\begin{description}
\item[{\em n}]Number of table entries \item[{\em r}]pointer to the red vector \item[{\em g}]pointer to the green vector \item[{\em b}]pointer to the blue vector \item[{\em flag}]0/1 for image/vector color load\end{description}
\end{Desc}
Loads the X-windows color palette given a set of RGB color vectors. \index{xtvutils.c@{xtvutils.c}!imageread@{imageread}}
\index{imageread@{imageread}!xtvutils.c@{xtvutils.c}}
\subsubsection{\setlength{\rightskip}{0pt plus 5cm}int imageread (int $\ast$ {\em x}, int $\ast$ {\em y}, char $\ast$ {\em key})}\label{xtvutils_8c_a188}


Blocking read from the display of one character. 

\begin{Desc}
\item[Parameters:]
\begin{description}
\item[{\em x}]User X-axis coordinate where the key was struck \item[{\em y}]User Y-axis coordinate where the key was struck \item[{\em key}]ASCII code of the key struck\end{description}
\end{Desc}
Puts the cursor on the image and waits for a key to be hit (used, e.g., to prompt for user cursor interaction on the image).

Blocking is rude, so do it sparingly. \index{xtvutils.c@{xtvutils.c}!imagerelocate@{imagerelocate}}
\index{imagerelocate@{imagerelocate}!xtvutils.c@{xtvutils.c}}
\subsubsection{\setlength{\rightskip}{0pt plus 5cm}void imagerelocate (int {\em xuser}, int {\em yuser})}\label{xtvutils_8c_a178}


Move the overlay vector pointer to user X,Y coordinates. 

\begin{Desc}
\item[Parameters:]
\begin{description}
\item[{\em xuser}]X coordinate in user space \item[{\em yuser}]Y coordinate in user space\end{description}
\end{Desc}
Moves the image overlay plane \char`\"{}pen\char`\"{} to a given (X,Y) location in user coordinates. A \char`\"{}relocate\char`\"{} specifies a starting point for a draw.

\begin{Desc}
\item[See also:]{\bf imagedraw()}{\rm (p.\,\pageref{xtv_8h_a73})} \end{Desc}
\index{xtvutils.c@{xtvutils.c}!imagetext@{imagetext}}
\index{imagetext@{imagetext}!xtvutils.c@{xtvutils.c}}
\subsubsection{\setlength{\rightskip}{0pt plus 5cm}void {\bf imagetext} (int {\em xuser}, int {\em yuser}, char $\ast$ {\em text}, int $\ast$ {\em textlen})}\label{xtvutils_8c_a175}


imagetext - draw text over the image 

\begin{Desc}
\item[Parameters:]
\begin{description}
\item[{\em xuser}]X-axis coordinate (user coords) of the start of the string \item[{\em yuser}]Y-axis coordinate (user coords) of the start of the string \item[{\em text}]string with the text \item[{\em textlen}]length of the string\end{description}
\end{Desc}
Draw text over an image. \index{xtvutils.c@{xtvutils.c}!imageupdate@{imageupdate}}
\index{imageupdate@{imageupdate}!xtvutils.c@{xtvutils.c}}
\subsubsection{\setlength{\rightskip}{0pt plus 5cm}int imageupdate (int {\em x0}, int {\em y0}, int {\em nx}, int {\em ny}, int {\em map})}\label{xtvutils_8c_a169}


update the image (or a subset) on the display 

\begin{Desc}
\item[Parameters:]
\begin{description}
\item[{\em x0}]starting X-axis pixel value of the region to update \item[{\em y0}]starting Y-axis pixel value of the region to update \item[{\em nx}]number of X-axis pixels in the region \item[{\em ny}]number of Y-axis pixels in the region \item[{\em map}]flag, if =1, means to remap data values into display bits\end{description}
\end{Desc}
\begin{Desc}
\item[Returns:]0 on success, error codes on failures.\end{Desc}
Updates all or part of an image on the display. \index{xtvutils.c@{xtvutils.c}!keyhelp@{keyhelp}}
\index{keyhelp@{keyhelp}!xtvutils.c@{xtvutils.c}}
\subsubsection{\setlength{\rightskip}{0pt plus 5cm}void keyhelp (int {\em x}, int {\em y}, int {\em xuser}, int {\em yuser}, int {\em key})}\label{xtvutils_8c_a215}


Help Action Function. 

\begin{Desc}
\item[Parameters:]
\begin{description}
\item[{\em x}]X-axis position of the cursor in window space \item[{\em y}]Y-axis position of the cursor in window space \item[{\em xuser}]X-axis position of the cursor in user space \item[{\em yuser}]Y-axis position of the cursor in user space \item[{\em key}]ASCII code of the key pressed\end{description}
\end{Desc}
Print the help menu for action keys \index{xtvutils.c@{xtvutils.c}!keypan@{keypan}}
\index{keypan@{keypan}!xtvutils.c@{xtvutils.c}}
\subsubsection{\setlength{\rightskip}{0pt plus 5cm}void keypan (int {\em x}, int {\em y}, int {\em xuser}, int {\em yuser}, int {\em key})}\label{xtvutils_8c_a212}


Image Pan Action Function. 

\begin{Desc}
\item[Parameters:]
\begin{description}
\item[{\em x}]X-axis position of the cursor in window space \item[{\em y}]Y-axis position of the cursor in window space \item[{\em xuser}]X-axis position of the cursor in user space \item[{\em yuser}]Y-axis position of the cursor in user space \item[{\em key}]ASCII code of the key pressed\end{description}
\end{Desc}
Pan the image to center on the cursor position \index{xtvutils.c@{xtvutils.c}!keyrecenter@{keyrecenter}}
\index{keyrecenter@{keyrecenter}!xtvutils.c@{xtvutils.c}}
\subsubsection{\setlength{\rightskip}{0pt plus 5cm}void keyrecenter (int {\em x}, int {\em y}, int {\em xuser}, int {\em yuser}, int {\em key})}\label{xtvutils_8c_a213}


Restore Original Zoom/Center Action Function. 

\begin{Desc}
\item[Parameters:]
\begin{description}
\item[{\em x}]X-axis position of the cursor in window space \item[{\em y}]Y-axis position of the cursor in window space \item[{\em xuser}]X-axis position of the cursor in user space \item[{\em yuser}]Y-axis position of the cursor in user space \item[{\em key}]ASCII code of the key pressed\end{description}
\end{Desc}
De-Zoom and Recenter the image in the window (aka \char`\"{}restore original zoom/pan on initial display). \index{xtvutils.c@{xtvutils.c}!keyzoomin@{keyzoomin}}
\index{keyzoomin@{keyzoomin}!xtvutils.c@{xtvutils.c}}
\subsubsection{\setlength{\rightskip}{0pt plus 5cm}void keyzoomin (int {\em x}, int {\em y}, int {\em xuser}, int {\em yuser}, int {\em key})}\label{xtvutils_8c_a210}


Zoom In Action Function. 

\begin{Desc}
\item[Parameters:]
\begin{description}
\item[{\em x}]X-axis position of the cursor in window space \item[{\em y}]Y-axis position of the cursor in window space \item[{\em xuser}]X-axis position of the cursor in user space \item[{\em yuser}]Y-axis position of the cursor in user space \item[{\em key}]ASCII code of the key pressed\end{description}
\end{Desc}
Action function to zoom in 1 step \index{xtvutils.c@{xtvutils.c}!keyzoomout@{keyzoomout}}
\index{keyzoomout@{keyzoomout}!xtvutils.c@{xtvutils.c}}
\subsubsection{\setlength{\rightskip}{0pt plus 5cm}void keyzoomout (int {\em x}, int {\em y}, int {\em xuser}, int {\em yuser}, int {\em key})}\label{xtvutils_8c_a211}


Zoom Out Action Function. 

\begin{Desc}
\item[Parameters:]
\begin{description}
\item[{\em x}]X-axis position of the cursor in window space \item[{\em y}]Y-axis position of the cursor in window space \item[{\em xuser}]X-axis position of the cursor in user space \item[{\em yuser}]Y-axis position of the cursor in user space \item[{\em key}]ASCII code of the key pressed\end{description}
\end{Desc}
Action function to zoom out 1 step \index{xtvutils.c@{xtvutils.c}!keyzoompan@{keyzoompan}}
\index{keyzoompan@{keyzoompan}!xtvutils.c@{xtvutils.c}}
\subsubsection{\setlength{\rightskip}{0pt plus 5cm}void keyzoompan (int {\em x}, int {\em y}, int {\em inout})}\label{xtvutils_8c_a209}


Zoom/Pan Key Action Function. 

\begin{Desc}
\item[Parameters:]
\begin{description}
\item[{\em x}]X-axis coordinates of the cursor in window space \item[{\em y}]Y-axis coordinates of the cursor in window space \item[{\em inout}]amount to zoom in (+1) or out (-1)\end{description}
\end{Desc}
Low-level zoom and pan function. \index{xtvutils.c@{xtvutils.c}!keyzoomprint@{keyzoomprint}}
\index{keyzoomprint@{keyzoomprint}!xtvutils.c@{xtvutils.c}}
\subsubsection{\setlength{\rightskip}{0pt plus 5cm}void keyzoomprint (int {\em x}, int {\em y}, int {\em xuser}, int {\em yuser}, int {\em key})}\label{xtvutils_8c_a214}


Power Zoom/Pan Action Function. 

\begin{Desc}
\item[Parameters:]
\begin{description}
\item[{\em x}]X-axis position of the cursor in window space \item[{\em y}]Y-axis position of the cursor in window space \item[{\em xuser}]X-axis position of the cursor in user space \item[{\em yuser}]Y-axis position of the cursor in user space \item[{\em key}]ASCII code of the key pressed\end{description}
\end{Desc}
Power Zoom/Pan - zooms right in on the pixel under the cursor, recentering on the display. \index{xtvutils.c@{xtvutils.c}!lastim@{lastim}}
\index{lastim@{lastim}!xtvutils.c@{xtvutils.c}}
\subsubsection{\setlength{\rightskip}{0pt plus 5cm}void lastim (int {\em x}, int {\em y}, int {\em xuser}, int {\em yuser}, int {\em key})}\label{xtvutils_8c_a218}


Blink to Previous Image Action Function. 

\begin{Desc}
\item[Parameters:]
\begin{description}
\item[{\em x}]X-axis position of the cursor in window space \item[{\em y}]Y-axis position of the cursor in window space \item[{\em xuser}]X-axis position of the cursor in user space \item[{\em yuser}]Y-axis position of the cursor in user space \item[{\em key}]ASCII code of the key pressed\end{description}
\end{Desc}
Blink back to the previous image in the ring buffer. \index{xtvutils.c@{xtvutils.c}!lights@{lights}}
\index{lights@{lights}!xtvutils.c@{xtvutils.c}}
\subsubsection{\setlength{\rightskip}{0pt plus 5cm}void lights (int {\em state})}\label{xtvutils_8c_a186}


Work the various status \char`\"{}lights\char`\"{} on the display. 

\begin{Desc}
\item[Parameters:]
\begin{description}
\item[{\em state}]status light to operate and its state: +=on/-=off\end{description}
\end{Desc}
The display status lights are numbered 1..4 running from top to bottom. In order: \small\begin{alltt}
  Light 1: Input State (-1=Asychronous, +1=waiting for key press)
  Light 2: Image State (-2=no image, 2=image loaded \& ready)
  Light 3: Zoom Window State (-3=frozen, +3=active)
  Light 4: Image Zoom State (-4 = normal, +4=zoomed in/out)
  \end{alltt}\normalsize 
 In the normal state the light is green, otherwise it is red.

If state=0, it turns all the lights on GREEN (+state) \index{xtvutils.c@{xtvutils.c}!mapimage@{mapimage}}
\index{mapimage@{mapimage}!xtvutils.c@{xtvutils.c}}
\subsubsection{\setlength{\rightskip}{0pt plus 5cm}void mapimage (int {\em x0}, int {\em y0}, int {\em nx}, int {\em ny}, int {\em datw}, float $\ast$ {\em data}, int {\em imx0}, int {\em imy0}, int {\em imw}, int {\em yinv}, unsigned long $\ast$ {\em image}, float $\ast$ {\em breakpts}, unsigned long {\em nbreak}, unsigned long $\ast$ {\em pixels}, int {\em bits\_\-per\_\-pixel}, unsigned long {\em cmask})}\label{xtvutils_8c_a224}


Map floating image data into integer display data. 

\begin{Desc}
\item[Parameters:]
\begin{description}
\item[{\em x0}]X-axis coordinate of origin of image section to map \item[{\em y0}]Y-axis coordinate of origin of image section to map \item[{\em nx}]X-axis size of image section to map \item[{\em ny}]Y-axis size of image section to map \item[{\em datw}]Number of data pixels per row in full image \item[{\em data}]array of data to be mapped \item[{\em imx0}]X-axis origin of the image array \item[{\em imy0}]Y-axis origin of the image array \item[{\em imw}]Number of image pixels per row \item[{\em yinv}]Flag: invert y order of data and image \item[{\em image}]array to receive mapped data \item[{\em breakpts}]lookup table \item[{\em nbreak}]fallback branch table for ambiguous cases \item[{\em pixels}]color cells vector (defined in zimage.h) \item[{\em bits\_\-per\_\-pixel}]data precision (bits per pixel), values: 8, 16, or 24 \item[{\em cmask}]color mask\end{description}
\end{Desc}
Uses the lookup table to map image data values into N-bit display values for display. Unlike an older version of {\bf mapimage()}{\rm (p.\,\pageref{xtv_8h_a119})}, this version knows about multiple color planes and color masks (e.g., for RGB color planes in a True\-Color visual like a 24-bit display device). \index{xtvutils.c@{xtvutils.c}!newcolors@{newcolors}}
\index{newcolors@{newcolors}!xtvutils.c@{xtvutils.c}}
\subsubsection{\setlength{\rightskip}{0pt plus 5cm}void newcolors (int {\em n1}, int {\em n2})}\label{xtvutils_8c_a201}


Compress the color palette. 

\begin{Desc}
\item[Parameters:]
\begin{description}
\item[{\em n1}]index of the new Min in the palette \item[{\em n2}]index of the new Max in the palette.\end{description}
\end{Desc}
Repacks the color palette (compresses it) to have the effect of changing the image intensity scaling (zeropoint and span), though at a sacrifice of bit depth, by re-mapping the intensities to only span palette indexes n1..n2 (normally it runs 0..npalette-1). \index{xtvutils.c@{xtvutils.c}!newsizesubwin@{newsizesubwin}}
\index{newsizesubwin@{newsizesubwin}!xtvutils.c@{xtvutils.c}}
\subsubsection{\setlength{\rightskip}{0pt plus 5cm}void newsizesubwin (int {\em wid}, int {\em hgt})}\label{xtvutils_8c_a193}


Resize subwindows. 

\begin{Desc}
\item[Parameters:]
\begin{description}
\item[{\em wid}]new subwindow width \item[{\em hgt}]new subwindow height \end{description}
\end{Desc}
\index{xtvutils.c@{xtvutils.c}!nextim@{nextim}}
\index{nextim@{nextim}!xtvutils.c@{xtvutils.c}}
\subsubsection{\setlength{\rightskip}{0pt plus 5cm}void nextim (int {\em x}, int {\em y}, int {\em xuser}, int {\em yuser}, int {\em key})}\label{xtvutils_8c_a217}


Blink to Next Image Action Function. 

\begin{Desc}
\item[Parameters:]
\begin{description}
\item[{\em x}]X-axis position of the cursor in window space \item[{\em y}]Y-axis position of the cursor in window space \item[{\em xuser}]X-axis position of the cursor in user space \item[{\em yuser}]Y-axis position of the cursor in user space \item[{\em key}]ASCII code of the key pressed\end{description}
\end{Desc}
Blink forward to the next image in the ring buffer. \index{xtvutils.c@{xtvutils.c}!oldzcursor@{oldzcursor}}
\index{oldzcursor@{oldzcursor}!xtvutils.c@{xtvutils.c}}
\subsubsection{\setlength{\rightskip}{0pt plus 5cm}void oldzcursor (int {\em j})}\label{xtvutils_8c_a205}


Old zoom-window cursor position (updatezoom() utility function). 

\begin{Desc}
\item[Parameters:]
\begin{description}
\item[{\em j}]the location of the top left of the central pixel \end{description}
\end{Desc}
\index{xtvutils.c@{xtvutils.c}!replicate@{replicate}}
\index{replicate@{replicate}!xtvutils.c@{xtvutils.c}}
\subsubsection{\setlength{\rightskip}{0pt plus 5cm}void replicate (int {\em x0}, int {\em y0}, int {\em w0}, int {\em h0}, char $\ast$ {\em from}, int {\em f}, int {\em x}, int {\em y}, int {\em w}, int {\em h}, int {\em wto}, char $\ast$ {\em to}, int {\em fill}, int {\em bits\_\-per\_\-pixel})}\label{xtvutils_8c_a225}


replicate 

\begin{Desc}
\item[Parameters:]
\begin{description}
\item[{\em x0}]X-axis origin of the source image \item[{\em y0}]Y-axis origin of the source image \item[{\em w0}]width (X-axis size) of source image \item[{\em h0}]height (Y-axis size) of source image \item[{\em from}]Source image array \item[{\em f}]replication factor \item[{\em x}]X-axis origin of the destination image \item[{\em y}]Y-axis origin of the destination image \item[{\em w}]width (X-axis size) of the destination image \item[{\em h}]height (Y-axis size) of the destination image \item[{\em wto}]dimensions of the destination image \item[{\em to}]Destination image array \item[{\em fill}]Flag: 1=zero rest of the to array \item[{\em bits\_\-per\_\-pixel}]data precision of the arrays\end{description}
\end{Desc}
Fills the destination array with an arbitrary (square) repetition of each source pixel into Nx\-N pixels in the destination array

Fill array \char`\"{}to\char`\"{} with a piece of \char`\"{}from\char`\"{}, each pixel mapping to fxf pixels If fill != 0, zero any uncovered piece of destination array

Note, this version uses memmove() instead of bcopy() since the latter is now officially deprecated. \index{xtvutils.c@{xtvutils.c}!resetcolors@{resetcolors}}
\index{resetcolors@{resetcolors}!xtvutils.c@{xtvutils.c}}
\subsubsection{\setlength{\rightskip}{0pt plus 5cm}void resetcolors ()}\label{xtvutils_8c_a200}


Reset the color palette. 

Restores the full color palette after compression

\begin{Desc}
\item[See also:]{\bf newcolors()}{\rm (p.\,\pageref{xtv_8h_a95})} \end{Desc}
\index{xtvutils.c@{xtvutils.c}!resizesubwin@{resizesubwin}}
\index{resizesubwin@{resizesubwin}!xtvutils.c@{xtvutils.c}}
\subsubsection{\setlength{\rightskip}{0pt plus 5cm}void resizesubwin ()}\label{xtvutils_8c_a194}


resize the subwindow 

Called by higher-level routines to do the actual dirty work \index{xtvutils.c@{xtvutils.c}!samplicate@{samplicate}}
\index{samplicate@{samplicate}!xtvutils.c@{xtvutils.c}}
\subsubsection{\setlength{\rightskip}{0pt plus 5cm}void samplicate (int {\em x0}, int {\em y0}, int {\em w0}, int {\em h0}, char $\ast$ {\em from}, int {\em f}, int {\em x}, int {\em y}, int {\em w}, int {\em h}, int {\em wto}, char $\ast$ {\em to}, int {\em fill}, int {\em bits\_\-per\_\-pixel})}\label{xtvutils_8c_a226}


samplicate an array (sample instead of replicate) 

\begin{Desc}
\item[Parameters:]
\begin{description}
\item[{\em x0}]X-axis origin of the source image \item[{\em y0}]Y-axis origin of the source image \item[{\em w0}]width (X-axis size) of source image \item[{\em h0}]height (Y-axis size) of source image \item[{\em from}]Source image array \item[{\em f}]sampling factor \item[{\em x}]X-axis origin of the destination image \item[{\em y}]Y-axis origin of the destination image \item[{\em w}]width (X-axis size) of the destination image \item[{\em h}]height (Y-axis size) of the destination image \item[{\em wto}]dimensions of the destination image \item[{\em to}]Destination image array \item[{\em fill}]Flag: 1=zero rest of the to array \item[{\em bits\_\-per\_\-pixel}]data precision of the arrays\end{description}
\end{Desc}
fills the destination array with an arbitrary (square) sampling sampling of each Nx\-N source pixels into the destination array \index{xtvutils.c@{xtvutils.c}!storeclear@{storeclear}}
\index{storeclear@{storeclear}!xtvutils.c@{xtvutils.c}}
\subsubsection{\setlength{\rightskip}{0pt plus 5cm}void storeclear ()}\label{xtvutils_8c_a182}


clears the image storage (ring) buffer 

Clears all images from the image ring buffer, used to store multiple images for blinking. \index{xtvutils.c@{xtvutils.c}!tvimnum@{tvimnum}}
\index{tvimnum@{tvimnum}!xtvutils.c@{xtvutils.c}}
\subsubsection{\setlength{\rightskip}{0pt plus 5cm}void tvimnum (char $\ast$ {\em str}, int {\em n})}\label{xtvutils_8c_a196}


Update the image number of the displayed image. 

\begin{Desc}
\item[Parameters:]
\begin{description}
\item[{\em str}]string with the image number \item[{\em n}]integer version of str\end{description}
\end{Desc}
XVista legacy routines, not implemented here but available if someday we want something like it. \index{xtvutils.c@{xtvutils.c}!updatebrkpt@{updatebrkpt}}
\index{updatebrkpt@{updatebrkpt}!xtvutils.c@{xtvutils.c}}
\subsubsection{\setlength{\rightskip}{0pt plus 5cm}void updatebrkpt (int {\em x}, int {\em ibut})}\label{xtvutils_8c_a198}


Display the color palette break point vector value on cursor action. 

\begin{Desc}
\item[Parameters:]
\begin{description}
\item[{\em x}]X-axis coorindate of the cursor in the color palette window \item[{\em ibut}]ASCII code of any button pressed \end{description}
\end{Desc}
\index{xtvutils.c@{xtvutils.c}!updatecoords@{updatecoords}}
\index{updatecoords@{updatecoords}!xtvutils.c@{xtvutils.c}}
\subsubsection{\setlength{\rightskip}{0pt plus 5cm}void updatecoords (int {\em wx}, int {\em wy}, int {\em key})}\label{xtvutils_8c_a197}


Update the cursor x,y coords on the status window. 

\begin{Desc}
\item[Parameters:]
\begin{description}
\item[{\em wx}]X-axis position of the cursor in window space \item[{\em wy}]Y-axis position of the cursor in window space \item[{\em key}]ASCII code of any action key hit\end{description}
\end{Desc}
Updates the X,Y coords and associated data in the status subwindow based on the cursor position and any key that might have been struck. Used as a utility function for any key event action functions that will update the contents of the status window. \index{xtvutils.c@{xtvutils.c}!updateimage@{updateimage}}
\index{updateimage@{updateimage}!xtvutils.c@{xtvutils.c}}
\subsubsection{\setlength{\rightskip}{0pt plus 5cm}void updateimage (int {\em x0}, int {\em y0}, int {\em nx}, int {\em ny}, int {\em awidth}, float $\ast$ {\em a}, int {\em x1}, int {\em y1})}\label{xtvutils_8c_a189}


Update the image and window dimension database. 

\begin{Desc}
\item[Parameters:]
\begin{description}
\item[{\em x0}]X-axis location of the image subsection in the data array \item[{\em y0}]X-axis location of the image subsection in the data array \item[{\em nx}]X-axis size of the image subsection in the data array \item[{\em ny}]Y-axis size of the image subsection in the data array \item[{\em awidth}]number of pixels/line in the data array \item[{\em a}]pointer to the floating-point data array \item[{\em x1}]X-axis offset of the data coordinate system \item[{\em y1}]Y-axis offset of the data coordinate system\end{description}
\end{Desc}
Updates the image and window dimension database after a resize or other operations that change the size of the window for new images. \index{xtvutils.c@{xtvutils.c}!updatename@{updatename}}
\index{updatename@{updatename}!xtvutils.c@{xtvutils.c}}
\subsubsection{\setlength{\rightskip}{0pt plus 5cm}void updatename (char $\ast$ {\em text}, int {\em num})}\label{xtvutils_8c_a195}


update the name that appears on the image display 

\begin{Desc}
\item[Parameters:]
\begin{description}
\item[{\em text}]string with the new image name \item[{\em num}]number of the text string to change (1 or 2)\end{description}
\end{Desc}
Changes one or other of the label strings for the xyz status subwindow (displays cursor x,y and corresponding data values). \small\begin{alltt}
  num=1 sets imname1[], the top label
  num=2 sets imname2[], the middle label
  \end{alltt}\normalsize 
 Any other values of num are ignored.

Note that the bottom label in the xzy subwindow is either the position/data readout if the cursor is on the image, or the color lookup table value if the cursor is on the color palette subwindow. \index{xtvutils.c@{xtvutils.c}!updatenewpal@{updatenewpal}}
\index{updatenewpal@{updatenewpal}!xtvutils.c@{xtvutils.c}}
\subsubsection{\setlength{\rightskip}{0pt plus 5cm}void updatenewpal (int {\em x0}, int {\em ibut})}\label{xtvutils_8c_a202}


update the color palette after modification 

\begin{Desc}
\item[Parameters:]
\begin{description}
\item[{\em x0}]position of breakpoint in the color palette \item[{\em ibut}]button pressed (?) \end{description}
\end{Desc}
\index{xtvutils.c@{xtvutils.c}!updatepal@{updatepal}}
\index{updatepal@{updatepal}!xtvutils.c@{xtvutils.c}}
\subsubsection{\setlength{\rightskip}{0pt plus 5cm}void updatepal (int {\em x}, int {\em ibut})}\label{xtvutils_8c_a203}


Update the palette. 

\begin{Desc}
\item[Parameters:]
\begin{description}
\item[{\em x}]repack the color palette to this breakpoint \item[{\em ibut}]action button hit (?) \end{description}
\end{Desc}
\index{xtvutils.c@{xtvutils.c}!updatepan@{updatepan}}
\index{updatepan@{updatepan}!xtvutils.c@{xtvutils.c}}
\subsubsection{\setlength{\rightskip}{0pt plus 5cm}void updatepan (int {\em xc}, int {\em yc}, int {\em in})}\label{xtvutils_8c_a191}


Update the centering of the image (pan-and-zoom). 

\begin{Desc}
\item[Parameters:]
\begin{description}
\item[{\em xc}]new image center in X \item[{\em yc}]new image center in Y \item[{\em in}]zoom: 1=zoom in, -1=zoom out\end{description}
\end{Desc}
Pans the center of the image to user coordinate (xc,yc) and either zooms in 1 step (in=1) or zooms out 1 step (in=-1). \index{xtvutils.c@{xtvutils.c}!updatesize@{updatesize}}
\index{updatesize@{updatesize}!xtvutils.c@{xtvutils.c}}
\subsubsection{\setlength{\rightskip}{0pt plus 5cm}void updatesize (int {\em wid}, int {\em hgt})}\label{xtvutils_8c_a192}


Update all the windows to a new overall size. 

\begin{Desc}
\item[Parameters:]
\begin{description}
\item[{\em wid}]new window width \item[{\em hgt}]new window height\end{description}
\end{Desc}
Sets the new width and height for the window and its associated subframes. \index{xtvutils.c@{xtvutils.c}!vecclear@{vecclear}}
\index{vecclear@{vecclear}!xtvutils.c@{xtvutils.c}}
\subsubsection{\setlength{\rightskip}{0pt plus 5cm}void vecclear ()}\label{xtvutils_8c_a181}


Clear the overlay plane. 

Erases all lines drawn in the image overlay plane (resets the overlay vector) \index{xtvutils.c@{xtvutils.c}!whitePixel@{whitePixel}}
\index{whitePixel@{whitePixel}!xtvutils.c@{xtvutils.c}}
\subsubsection{\setlength{\rightskip}{0pt plus 5cm}int white\-Pixel (Display $\ast$ {\em dpy}, int {\em screen})}\label{xtvutils_8c_a223}


Border pixmap (white). 

\begin{Desc}
\item[Parameters:]
\begin{description}
\item[{\em dpy}]Pointer to the display \item[{\em screen}]screen number \end{description}
\end{Desc}
\index{xtvutils.c@{xtvutils.c}!writepix@{writepix}}
\index{writepix@{writepix}!xtvutils.c@{xtvutils.c}}
\subsubsection{\setlength{\rightskip}{0pt plus 5cm}void writepix (int {\em x}, int {\em y}, int {\em wid}, int {\em hgt})}\label{xtvutils_8c_a207}


Write a portion of the image to the image window. 

\begin{Desc}
\item[Parameters:]
\begin{description}
\item[{\em x}]X-axis coordinate of the region to display in window space \item[{\em y}]Y-axis coordinate of the region to display in window space \item[{\em wid}]X-axis width in window space \item[{\em hgt}]Y-axis height in window space\end{description}
\end{Desc}
A low-level (window space) image subsection display function, usually called only by functions that do all the conversion from user space.

This is where the pixels hit the screen. \index{xtvutils.c@{xtvutils.c}!xtv_refresh@{xtv\_\-refresh}}
\index{xtv_refresh@{xtv\_\-refresh}!xtvutils.c@{xtvutils.c}}
\subsubsection{\setlength{\rightskip}{0pt plus 5cm}int xtv\_\-refresh (int {\em signo})}\label{xtvutils_8c_a208}


Interrupt handler for X Events. 

\begin{Desc}
\item[Parameters:]
\begin{description}
\item[{\em signo}]integer signal code to handle\end{description}
\end{Desc}
This is the function that actually handles actions on window events (cursor moving over image, buttons getting pressed, etc.) \index{xtvutils.c@{xtvutils.c}!xy2index@{xy2index}}
\index{xy2index@{xy2index}!xtvutils.c@{xtvutils.c}}
\subsubsection{\setlength{\rightskip}{0pt plus 5cm}int xy2index (int {\em xuser}, int {\em yuser})}\label{xtvutils_8c_a231}


Convert (x,y) user coordinates into data vector index. 

\begin{Desc}
\item[Parameters:]
\begin{description}
\item[{\em xuser}]user X coordinate \item[{\em yuser}]user Y coordinate \end{description}
\end{Desc}
\begin{Desc}
\item[Returns:]array index corresponding to (x,y), or -1 if out of bounds \end{Desc}
\index{xtvutils.c@{xtvutils.c}!zcursor@{zcursor}}
\index{zcursor@{zcursor}!xtvutils.c@{xtvutils.c}}
\subsubsection{\setlength{\rightskip}{0pt plus 5cm}void zcursor (int {\em j})}\label{xtvutils_8c_a206}


New zoom cursor position (updatezoom() utility function). 

\begin{Desc}
\item[Parameters:]
\begin{description}
\item[{\em j}]the location of the top left of the central pixel \end{description}
\end{Desc}
\index{xtvutils.c@{xtvutils.c}!zeroborder@{zeroborder}}
\index{zeroborder@{zeroborder}!xtvutils.c@{xtvutils.c}}
\subsubsection{\setlength{\rightskip}{0pt plus 5cm}void zeroborder (int {\em x0}, int {\em y0}, int {\em w0}, int {\em h0}, int {\em w}, int {\em h}, int {\em wto}, char $\ast$ {\em to}, int {\em nbytes})}\label{xtvutils_8c_a228}


fill border with zeros 

\begin{Desc}
\item[Parameters:]
\begin{description}
\item[{\em x0}]X-axis origin of the source image \item[{\em y0}]Y-axis origin of the source image \item[{\em w0}]width (X-axis size) of source image \item[{\em h0}]height (Y-axis size) of source image \item[{\em w}]width (X-axis size) of the destination image \item[{\em h}]height (Y-axis size) of the destination image \item[{\em wto}]dimensions of the destination image \item[{\em to}]Destination image array \item[{\em nbytes}]\end{description}
\end{Desc}
\index{xtvutils.c@{xtvutils.c}!zfreeze@{zfreeze}}
\index{zfreeze@{zfreeze}!xtvutils.c@{xtvutils.c}}
\subsubsection{\setlength{\rightskip}{0pt plus 5cm}void zfreeze (int {\em x}, int {\em y}, int {\em xuser}, int {\em yuser}, int {\em key})}\label{xtvutils_8c_a216}


Freeze Zoom Action Function. 

\begin{Desc}
\item[Parameters:]
\begin{description}
\item[{\em x}]X-axis position of the cursor in window space \item[{\em y}]Y-axis position of the cursor in window space \item[{\em xuser}]X-axis position of the cursor in user space \item[{\em yuser}]Y-axis position of the cursor in user space \item[{\em key}]ASCII code of the key pressed\end{description}
\end{Desc}
Routine that freezes the zoom window on one location \index{xtvutils.c@{xtvutils.c}!zhairs@{zhairs}}
\index{zhairs@{zhairs}!xtvutils.c@{xtvutils.c}}
\subsubsection{\setlength{\rightskip}{0pt plus 5cm}void zhairs (int {\em x}, int {\em y}, int {\em xuser}, int {\em yuser}, int {\em key})}\label{xtvutils_8c_a220}


Toggle Full-Screen Cross Hairs On/Off Action Function. 

\begin{Desc}
\item[Parameters:]
\begin{description}
\item[{\em x}]X-axis position of the cursor in window space \item[{\em y}]Y-axis position of the cursor in window space \item[{\em xuser}]X-axis position of the cursor in user space \item[{\em yuser}]Y-axis position of the cursor in user space \item[{\em key}]ASCII code of the key pressed\end{description}
\end{Desc}
Toggles between the little cross and full-screen cross hairs. \index{xtvutils.c@{xtvutils.c}!zpeak@{zpeak}}
\index{zpeak@{zpeak}!xtvutils.c@{xtvutils.c}}
\subsubsection{\setlength{\rightskip}{0pt plus 5cm}void zpeak (int {\em x}, int {\em y}, int {\em xuser}, int {\em yuser}, int {\em key})}\label{xtvutils_8c_a219}


Find Peak Pixel Nearest the Cursor Action Function. 

\begin{Desc}
\item[Parameters:]
\begin{description}
\item[{\em x}]X-axis position of the cursor in window space \item[{\em y}]Y-axis position of the cursor in window space \item[{\em xuser}]X-axis position of the cursor in user space \item[{\em yuser}]Y-axis position of the cursor in user space \item[{\em key}]ASCII code of the key pressed\end{description}
\end{Desc}
Finds the peak pixel nearest the cursor and puts the cursor on that pixel. Searches +/-7 pixels around the current location. 